\documentclass{ximera}
% xmPreamble.tex en xmPrintstyle.sty worden automatisch geladen door ximera.cls (v2.7.8);
% NIET nog eens \input{./preamble.tex} of \addPrintStyle{.} doen (geeft 'already defined'-fouten).

\makeatletter % eigen img-map toevoegen aan het zoekpad (werkt standalone én vanuit de xourse)
\g@addto@macro\Ginput@path{{img/}{UW4202/img/}}
\makeatother

\begin{document}
\author{Alexander Holvoet}
\xmtitle{Zand op een cirkelsegment}{}

\begin{image}[5cm]
    \includegraphics[width=5cm,keepaspectratio]{zand16.jpg}
\end{image}
Een cirkelsegment als grondvlak.

We nemen als basis voor een zandhoop een deel van een cirkelschijf aan \'e\'en kant van een
koorde. Deze figuur noemen we een \emph{cirkelsegment}. Het weggesneden stuk is ook een
cirkelsegment. We nemen hier een `groot' cirkelsegment \textemdash{} dit wil zeggen het
stuk met het middelpunt erop \textemdash{} maar het gaat even goed met een `klein'
cirkelsegment.

\begin{exercise}
    Teken (het bovenaanzicht van) de kamlijn die je verwacht, giet dan zand op het
    kartonnen exemplaar om dit te verifi\"eren, en bepaal waaraan de punten van de kamlijn
    voldoen.
    \begin{oplossing}
        Voor elk punt \(P\) van de kamlijn is de afstand \(|PQ|\) tot de koorde gelijk aan
        de afstand \(|PR|\) tot de rand van de cirkel.
        \begin{image}[5cm]
            \includegraphics[width=5cm,keepaspectratio]{zand17.jpg}
        \end{image}
    \end{oplossing}
\end{exercise}

Je vermoedt waarschijnlijk dat (het bovenaanzicht van) de kamlijn een parabool is. In de
volgende opgaven bewijs je dit.

\begin{exercise}
    Om te bewijzen dat de kamlijn een parabool is, heb je een brandpunt en een richtlijn
    nodig. Het brandpunt is het middelpunt \(M\) van de cirkel. Waar ongeveer moet de
    richtlijn \(d\) liggen? Toon dit met je vinger.
    \begin{oplossing}
        De richtlijn is een rechte evenwijdig met de koorde (vanwege de symmetrie) en
        `onder' de kamlijn. De precieze plaats komt in de volgende vraag aan bod.
    \end{oplossing}
\end{exercise}

\begin{exercise}
    Bewijs, door een geschikte keuze van de richtlijn, dat de kamlijn inderdaad een
    parabool is.
    \begin{hint}
        Neem een willekeurig punt \(P\) van de kamlijn en noem \(r\) de straal. Gebruik dat
        \(|PR| = |PQ|\), dus \(|PM| = r - |PR| = r - |PQ|\). Op welke afstand van de koorde
        moet je \(d\) leggen opdat \(|PM| = |PS|\)?
    \end{hint}
    \begin{oplossing}
        Neem een willekeurig punt \(P\) van de kamlijn en noem \(r\) de straal van de
        cirkel. We bepalen de richtlijn \(d\) zo dat de afstand \(|PM|\) tot het brandpunt
        gelijk is aan de afstand \(|PS|\) tot de rechte \(d\). We weten dat \(|PR| = |PQ|\),
        en dus
        \[
            |PM| = r - |PR| = r - |PQ| .
        \]
        Om \(|PM| = |PS|\) te krijgen, moeten we \(d\) op afstand \(r\) van de koorde nemen.
        Dan geldt inderdaad
        \[
            |PM| = r - |PR| = r - |PQ| = |QS| - |PQ| = |PS| .
        \]
        \begin{image}[5cm]
            \includegraphics[width=5cm,keepaspectratio]{zand18.jpg}
        \end{image}
    \end{oplossing}
\end{exercise}

\begin{exercise}
    Construeer de kamlijn met hoogtelijnen.
    \begin{image}[5cm]
        \includegraphics[width=5cm,keepaspectratio]{zand19.jpg}
    \end{image}
    \begin{oplossing}
        Construeer de snijpunten van de overeenkomstige hoogtelijnen. De rechte hoogtelijnen
        horen bij een schuin vlak dat volgens de rusthoek afloopt naar de koorde van het
        grondvlak. De cirkelvormige hoogtelijnen horen bij de kegel die je zou hebben als de
        basis de volledige cirkelschijf was. De kamlijn is de doorsnede van die kegel met
        het schuine vlak (even schuin als de beschrijvenden van de kegel). Dit is nog een
        extra manier om in te zien dat het een stuk parabool is.
        \begin{image}[15cm]
            \includegraphics[width=5cm,keepaspectratio]{zand20.jpg}
            \quad
            \includegraphics[width=5cm,keepaspectratio]{zand21.jpg}
            \quad
            \includegraphics[width=5cm,keepaspectratio]{zand22.jpg}
        \end{image}
    \end{oplossing}
\end{exercise}

\end{document}
